\documentclass[11pt]{amsart}

\usepackage[T1]{fontenc}
\usepackage{lmodern}
\usepackage{microtype}
\usepackage{amsmath,amssymb,mathtools}
\usepackage[hidelinks]{hyperref}

\hypersetup{
  pdftitle={A counterexample to Taranenko's diagonal-decomposition conjecture},
  pdfkeywords={polystochastic matrix, diagonal decomposition, Birkhoff polytope}
}

\newtheorem{theorem}{Theorem}
\newtheorem{proposition}[theorem]{Proposition}
\theoremstyle{remark}
\newtheorem{remark}[theorem]{Remark}

\DeclareMathOperator{\supp}{supp}

\title[Taranenko's diagonal-decomposition conjecture]
{A counterexample to Taranenko's diagonal-decomposition conjecture}
\date{}
\subjclass[2020]{15B51, 15A15, 52B05}
\keywords{polystochastic matrix, diagonal decomposition, Birkhoff polytope}

\begin{document}

\begin{abstract}
Taranenko conjectured that every polystochastic matrix of even dimension is a
nonnegative linear combination of diagonal indicators. We show that the
symmetric vertex $A_4\in\Omega_3^4$ from Taranenko's 2025 enumeration is not
in this cone. The proof is an explicit linear functional that is positive on
$A_4$ and nonpositive on every diagonal indicator. This disproves
Taranenko's Conjecture~4, restated by Child and Wanless as Conjecture~1.9.
The pair $(d,n)=(4,3)$ is componentwise minimal.
\end{abstract}

\maketitle

Let $I_n=\{0,\ldots,n-1\}$. A $d$-dimensional matrix of order $n$ is an
array $X=(X_x)_{x\in I_n^d}$. It is \emph{polystochastic} if $X_x\geq0$ and
every line has sum $1$; write $\Omega_n^d$ for the set of such arrays. A
\emph{diagonal} is a set $D\subset I_n^d$ of size $n$ such that every
coordinate projection maps $D$ bijectively onto $I_n$. Let $P_D$ be its
indicator array and set
\[
  \mathcal C_n^d=\operatorname{cone}\{P_D:D\text{ is a diagonal in }I_n^d\}.
\]
The arrays $P_D$ are the $(d-1)$-permutation matrices in the terminology of
Child and Wanless \cite{ChildWanless2020}.

Taranenko conjectured that $\Omega_n^d\subseteq\mathcal C_n^d$ whenever
$d$ is even \cite[Conjecture~4]{Taranenko2016}. Child and Wanless restated
this as Conjecture~1.9 in \cite{ChildWanless2020}; it also appears as
Conjecture~16 in Taranenko's 2025 dissertation
\cite{TaranenkoDissertation}. We give a counterexample in $\Omega_3^4$.

Put $I=I_3$. For $x\in I^4$, define
\[
  \nu(x)=(\nu_0(x),\nu_1(x),\nu_2(x)),
  \qquad
  \nu_r(x)=\bigl|\{j:x_j=r\}\bigr|.
\]
Let $A=(A_x)_{x\in I^4}$ be given by
\begin{equation}\label{eq:A}
A_x=
\begin{cases}
1,
  &\nu(x)\in\{(4,0,0),(0,4,0)\},\\[1mm]
\tfrac12,
  &\nu(x)\in\mathcal H,\\[1mm]
0,
  &\text{otherwise},
\end{cases}
\end{equation}
where
\begin{equation}\label{eq:H}
\mathcal H=
\bigl\{
(0,1,3),(0,2,2),(1,0,3),(1,2,1),
(2,0,2),(2,1,1),(2,2,0)
\bigr\}.
\end{equation}
With rows indexed lexicographically by $(x_1,x_2)$ and columns by
$(x_3,x_4)$, this is the symmetric vertex $A_4$ in
\cite[Appendix, item~4]{Taranenko2025}.

\begin{theorem}\label{thm:main}
The array $A$ belongs to $\Omega_3^4\setminus\mathcal C_3^4$. Consequently,
Conjecture~4 of \cite{Taranenko2016} and Conjecture~1.9 of
\cite{ChildWanless2020} are false.
\end{theorem}

\begin{proof}
Fix three coordinates and let $\mu=(\mu_0,\mu_1,\mu_2)$ record their
multiplicities. If $e_r$ is the $r$th standard basis vector, then the three
cells in the resulting line have types $\mu+e_0$, $\mu+e_1$, and
$\mu+e_2$. For $\mu=(3,0,0)$ and $(0,3,0)$, the corresponding values are
respectively $(1,0,0)$ and $(0,1,0)$. For each of the other eight weak
compositions of $3$, they are a permutation of
$(\tfrac12,\tfrac12,0)$. Thus every line has sum $1$, and
$A\in\Omega_3^4$.

Let $S=\supp(A)$ and define
\[
  S_+=\{x:\nu(x)\in\{(0,1,3),(1,0,3)\}\},
  \qquad
  S_-=\{x:\nu(x)=(2,2,0)\}.
\]
For any array $X$, set
\begin{equation}\label{eq:g}
  g(X)=
  \sum_{x\in S_+}X_x
  -\sum_{x\in S_-}X_x
  -\sum_{x\notin S}X_x.
\end{equation}
Since $|S_+|=2\binom41=8$ and $|S_-|=\binom42=6$, all at value
$\tfrac12$ in $A$,
\begin{equation}\label{eq:gA}
  g(A)=8\cdot\tfrac12-6\cdot\tfrac12=1.
\end{equation}

We show that $g(P_D)\leq0$ for every diagonal $D$. Every cell in $S_+$ has
three coordinates equal to $2$. Hence any two cells in $S_+$ agree in at
least one coordinate, so a diagonal contains at most one of them. It follows
immediately that $g(P_D)\leq0$ if $D\cap S_+=\varnothing$. The same holds if
$D\cap S_+\neq\varnothing$ but $D\nsubseteq S$: the possible contribution
$+1$ from $S_+$ is offset by at least one contribution $-1$ from
$I^4\setminus S$.

It remains to consider a diagonal $D\subseteq S$ meeting $S_+$. The supported
types are
\[
\begin{aligned}
\mathcal T=\{&(4,0,0),(0,4,0),(0,1,3),(0,2,2),(1,0,3),\\
              &(1,2,1),(2,0,2),(2,1,1),(2,2,0)\}.
\end{aligned}
\]
The three type vectors of the cells of $D$ sum to $(4,4,4)$, because each
coordinate position contains each symbol once along $D$. After fixing the
$S_+$-cell, the $\nu_2$-components of the other two types sum to $1$.
Among the supported types, those with $\nu_2=0$ are
\[
  (4,0,0),\quad(0,4,0),\quad(2,2,0),
\]
and those with $\nu_2=1$ are
\[
  (1,2,1),\quad(2,1,1).
\]
Matching the $\nu_0$- and $\nu_1$-components shows that the remaining types are
\[
\begin{cases}
(2,1,1)\text{ and }(2,2,0),&\text{if the $S_+$-type is }(0,1,3),\\
(1,2,1)\text{ and }(2,2,0),&\text{if the $S_+$-type is }(1,0,3).
\end{cases}
\]
Thus every supported diagonal meeting $S_+$ also meets $S_-$, and again
$g(P_D)\leq0$.

If $A=\sum_D c_D P_D$ with $c_D\geq0$, then linearity and
\eqref{eq:gA} give
\[
  1=g(A)=\sum_D c_D g(P_D)\leq0,
\]
a contradiction.
\end{proof}

\begin{proposition}\label{prop:minimal}
Within the class of even-dimensional polystochastic matrices, no
counterexample exists when $d=2$ or $n\leq2$. Hence $(d,n)=(4,3)$ is
componentwise minimal.
\end{proposition}

\begin{proof}
For $d=2$, the assertion is the Birkhoff--von Neumann theorem. Order $1$ is
trivial. Now let $X$ have order $2$ and even dimension $d$. Index its cells
by $x\in\{0,1\}^d$ and put $t=X_{(0,\ldots,0)}\in[0,1]$. Since the
two entries in every line sum to $1$, flipping one coordinate replaces a
cell value by its complement. Hence
\[
  X_x=
  \begin{cases}
    t,&\sum_i x_i\equiv0\pmod2,\\
    1-t,&\sum_i x_i\equiv1\pmod2.
  \end{cases}
\]
For $\bar x=(1-x_1,\ldots,1-x_d)$, the pair $\{x,\bar x\}$ is a diagonal.
Because $d$ is even, $x$ and $\bar x$ have the same parity. The complementary
pairs partition each parity class, and therefore
\[
  X=
  t\sum_{\substack{\{x,\bar x\}\text{ of even parity}}}
       P_{\{x,\bar x\}}
  +(1-t)\sum_{\substack{\{x,\bar x\}\text{ of odd parity}}}
       P_{\{x,\bar x\}},
\]
where each unordered pair is counted once.
\end{proof}

\begin{remark}[Scope and provenance]
The conjecture imposes no parity condition on the order, so the odd order of
$A$ is permitted. This example does not address the additional variant in
which the order is also required to be even. The array $A$ itself is not new;
the contribution here is the separating certificate \eqref{eq:g}.
Taranenko records $\operatorname{per}(A)=33/4>0$
\cite[Appendix, item~4]{Taranenko2025}, so the example does not concern the
separate positive-permanent conjecture.
\end{remark}

\begin{thebibliography}{9}

\bibitem{ChildWanless2020}
B.~Child and I.~M. Wanless,
\emph{Multidimensional permanents of polystochastic matrices},
Linear Algebra Appl. \textbf{586} (2020), 89--102.
\href{https://doi.org/10.1016/j.laa.2019.10.008}
{doi:10.1016/j.laa.2019.10.008}.

\bibitem{Taranenko2016}
A.~A. Taranenko,
\emph{Permanents of multidimensional matrices: properties and applications},
J. Appl. Ind. Math. \textbf{10} (2016), no.~4, 567--604.
\href{https://doi.org/10.1134/S1990478916040141}
{doi:10.1134/S1990478916040141}.

\bibitem{TaranenkoDissertation}
A.~A. Taranenko,
\emph{Multidimensional matrices in the algebraic theory of hypergraphs},
doctoral dissertation, Sobolev Institute of Mathematics, Novosibirsk, 2025
(in Russian).

\bibitem{Taranenko2025}
A.~A. Taranenko,
\emph{Enumeration and constructions of vertices of the polytope of
polystochastic matrices},
\href{https://arxiv.org/abs/2502.09149}
{arXiv:2502.09149 [math.CO]}, 2025.

\end{thebibliography}

\end{document}